\newpage

\section{Optimisation opérationnelle et superadditivité des économies}

Après avoir défini les décisions admissibles sur \(H\), nous déterminons le
coût minimal que peut atteindre chaque coalition. La répartition gloutonne
est d'abord établie pour une heure arbitraire ; elle devient ensuite la brique
élémentaire du problème persistant, dans lequel les gardiens restent fixes
sur \(H\). La politique horairement reconfigurable est traitée séparément.

L'optimisation comporte deux niveaux. Pour un ensemble de gardiens fixé, la
linéarité des coûts variables permet de déterminer explicitement la meilleure
répartition du trafic. Il reste ensuite à sélectionner les gardiens en tenant
compte de leurs coûts fixes de fonctionnement.

\subsection{Répartition optimale du trafic à gardiens fixés}
\label{subsec:allocation_gardiens_fixes}

Soit $(S, G)\in\mathcal E$. Les coûts fixes $\sum_{i\in G}F_i$ ne dépendent pas de la répartition du trafic, ce qui implique que minimiser le coût associé à la configuration de gardiens $(S, G)$ revient à minimiser ses coûts variables.

La répartition optimale du trafic pour $(S,G)$ est ainsi solution du
programme linéaire suivant :
\begin{equation}
\label{eq:lp_allocation_fixe}
\begin{aligned}
    \min_{\mathbf{t}\in\mathbb R_+^N}\quad
        & \sum_{i\in G}\gamma_i t_i,\\
    \text{s.c.}\quad
        & \sum_{i\in G} t_i=d(S),\\
        & 0\leqslant  t_i\leqslant q_i,
        \qquad \forall i\in G,\\
        & t_i=0,
        \qquad \forall i\in N\setminus G
\end{aligned}
\end{equation}


\subsubsection{Répartition gloutonne}

L'algorithme glouton répartit le trafic entre les équipements actifs par
ordre croissant de leur coût variable unitaire $\gamma_i$.

Soit $G=\{i_1,\dots,i_p\}$, ordonné de sorte que
$\gamma_{i_1}\leqslant\dots\leqslant\gamma_{i_p}$. En présence d'ex æquo,
une règle de départage déterministe est fixée. L'indice pivot est défini par
\[
k:=\min\left\{
r\in\{1,\dots,p\}\ \middle|\
\sum_{\ell=1}^{r}q_{i_\ell}\geqslant d(S)
\right\}.
\]
Cet ensemble est non vide puisque $G\in\mathcal G(S)$. La répartition
gloutonne $\mathbf t^g(S,G)$ sature les gardiens précédant le pivot,
affecte le trafic résiduel au pivot et laisse nul le trafic des gardiens
suivants :
\[
 t_{i_\ell}^g(S,G) =
\begin{cases} 
q_{i_\ell} & \text{si } \ell < k \\ 
d(S) - \sum_{m=1}^{k-1} q_{i_m} & \text{si } \ell = k \\
0 & \text{si } \ell > k
\end{cases}
\]
(\(t_i^g(S,G)=0\) pour tout \(i\notin G\)).
Lorsque \((S,G)\) est fixé dans un raisonnement, cette notation est abrégée
en \(\mathbf t^g\).

\subsubsection{Proposition d'optimalité}

\begin{proposition}[Optimalité de la répartition gloutonne]
\label{pro:opt}
    Soit $(S, G) \in \mathcal{E}$. Alors la répartition gloutonne du trafic
    $\mathbf{t}^g(S,G)$ minimise le coût opérationnel :
    \[
    C\bigl(S,G,\mathbf{t}^g(S,G)\bigr)
    =\min_{\mathbf{t}\in\mathcal T(S,G)}C(S,G,\mathbf{t})
    \]
\end{proposition}

\begin{proof}
Fixons \((S,G)\in\mathcal E\) et posons
\(\mathbf t^g:=\mathbf t^g(S,G)\). Soit
\(\mathbf t\in\mathcal T(S,G)\) une répartition admissible quelconque.

Les coûts fixes étant identiques pour les deux répartitions,
\[
\Delta
:=C(S,G,\mathbf t)-C(S,G,\mathbf t^g)
=\sum_{\ell=1}^{p}\gamma_{i_\ell}
\bigl(t_{i_\ell}-t_{i_\ell}^g\bigr).
\]
Comme les deux répartitions acheminent exactement \(d(S)\),
\(\sum_{\ell=1}^{p}(t_{i_\ell}-t_{i_\ell}^g)=0\). Si \(k\) désigne
l'indice pivot, il vient donc
\[
\Delta
=\sum_{\ell=1}^{p}
\bigl(\gamma_{i_\ell}-\gamma_{i_k}\bigr)
\bigl(t_{i_\ell}-t_{i_\ell}^g\bigr).
\]
Pour \(\ell<k\), les deux facteurs sont négatifs ou nuls, puisque
\(t_{i_\ell}^g=q_{i_\ell}\) ; pour \(\ell>k\), ils sont positifs ou nuls,
puisque \(t_{i_\ell}^g=0\). Le terme d'indice \(k\) est nul. Chaque terme
de la somme est donc positif ou nul, d'où \(\Delta\geqslant0\).
\end{proof}

La concentration du trafic résulte de la linéarité des coûts variables.
Si les $\gamma_i$ sont deux à deux distincts, la répartition optimale est
unique ; en présence d'ex æquo, plusieurs répartitions optimales peuvent
exister. Ce résultat vaut pour un ensemble de gardiens fixé et ne préjuge
ni de leur sélection, étudiée ci-dessous, ni de la répartition des économies
de coopération entre opérateurs.

En particulier, si tous les gardiens de \(G\) ont le même coût variable
unitaire \(\gamma\), l'objectif variable vaut \(\gamma d(S)\) pour toute
répartition admissible : chaque élément de \(\mathcal T(S,G)\) est donc
optimal. La règle de départage déterministe sélectionne seulement un
représentant opérationnel, sans modifier le coût minimal de la coalition.

Une fois les gardiens ordonnés, un parcours de la liste construit la
répartition gloutonne en $O(|G|)$. Si cet ordre doit être calculé, le tri
coûte en outre $O(|G|\log |G|)$ ; un tri global préalable des opérateurs
coûte $O(n\log n)$.







\subsection{Sélection persistante des gardiens sur la fenêtre}
\label{subsec:selection_gardiens}
\label{subsec:optimisation_fenetre}

Fixons \(S\in\mathcal S\). La sous-section précédente fournit la répartition
optimale du trafic à chaque heure lorsque l'ensemble de gardiens est fixé. Il
reste à choisir un même ensemble de gardiens pour toute la fenêtre.

Pour chaque opérateur $i\in S$, on introduit un indicateur binaire
d'activité $x_i$, égal à $1$ si son équipement est actif et à $0$ s'il
est en veille. Pour chaque \(h\in H\), la répartition du trafic est
représentée par \(\mathbf t^h\in\mathbb R_+^N\). Le vecteur
\(\mathbf x\) est commun à toutes les heures.

Le coût persistant optimal de la coalition $S$ est la valeur du programme
linéaire en nombres entiers mixtes suivant :
\begin{equation}
\label{eq:milp_gardiens}
\begin{aligned}
C_H^*(S)
=
\min_{\substack{\mathbf x\in\{0,1\}^S\\
                \mathbf t^h\in\mathbb R_+^N,\ h\in H}}\quad
    & \sum_{h\in H}\sum_{i\in S}\gamma_i t_i^h
    + |H|\sum_{i\in S}F_i x_i,\\
\text{s.c.}\quad
    & \sum_{i\in S}t_i^h=d^h(S),
      \qquad \forall h\in H,\\
    & 0\leqslant t_i^h\leqslant q_i x_i,
      \qquad \forall i\in S,\ \forall h\in H,\\
    & t_i^h=0,
      \qquad \forall i\in N\setminus S,\ \forall h\in H.
\end{aligned}
\end{equation}

La formulation \eqref{eq:milp_gardiens} est exacte : toute solution
réalisable définit
\(G=\{i\in S\mid x_i=1\}\in\mathcal G_H(S)\) et une répartition de
\(\mathcal T_h(S,G)\) à chaque heure. Réciproquement, toute solution
persistante admissible fournit une solution du programme en posant
\(x_i=\mathds{1}_{i\in G}\). Les fonctions objectif coïncident.



\subsection{Complexité et algorithme exact}

La présence des coûts fixes de fonctionnement et des indicateurs binaires d'activité rend le problème
combinatoire.

\begin{proposition}[NP-difficulté]
\label{prop:np_difficulte_gardiens}
Le problème de sélection persistante des gardiens est NP-difficile, même
lorsque tous les coûts variables unitaires sont identiques.
\end{proposition}

\emph{Démonstration : voir l'Annexe~\ref{app:preuve_np_difficulte}.}

\begin{remarque}
\textsc{Subset Sum} étant faiblement NP-complet, cette réduction établit
la NP-difficulté du problème, mais pas sa NP-difficulté forte.
\end{remarque}

Dans le cadre étudié, le nombre d'opérateurs présents sur un site est très
faible. Nous retenons donc une méthode exacte par énumération.

Pour chaque \(G\in\mathcal S\) tel que \(G\subseteq S\) :
\begin{enumerate}
    \item on teste sa faisabilité en vérifiant
    \(q(G)\geqslant d^h(S)\) pour tout \(h\in H\) ;
    \item si oui, on calcule la répartition gloutonne
    \(\mathbf t^{g,h}(S,G)\) à chaque heure ;
    \item on évalue son coût total
    \(\sum_{h\in H}C_h(S,G,\mathbf t^{g,h}(S,G))\) ;
    \item on conserve un ensemble de gardiens minimisant le coût.
\end{enumerate}

On obtient ainsi $C_H^*(S)$. Dans la suite, $G_H^*(S)$ désigne un
représentant arbitraire de
\[
\operatorname*{arg\,min}_{G\in\mathcal G_H(S)}
\sum_{h\in H}C_h\bigl(S,G,\mathbf t^{g,h}(S,G)\bigr),
\]
sans présumer son unicité, et
\(\mathbf t^{h,*}(S):=\mathbf t^{g,h}(S,G_H^*(S))\) désigne la répartition
associée à l'heure \(h\).

\begin{proposition}[Complexité de l'énumération]
\label{prop:complexite_enumeration}
Pour $S \in \mathcal S$ de taille $s$, l'énumération exacte des ensembles
de gardiens requiert \(O(|H|\,s\,2^s)\) opérations arithmétiques et
comparaisons. Le calcul des valeurs $C_H^*(S)$ pour toutes les
$S \in \mathcal S$ requiert \(O(|H|\,n\,3^n)\) telles opérations.
\end{proposition}

\emph{Démonstration : voir l'Annexe~\ref{app:preuve_complexite_enumeration}.}

\begin{remarque}
Le nombre de couples candidats $(S,G)$ est $\sum_{s=1}^{n}\binom ns(2^s - 1) = 3^n - 2^n.$
Cette quantité compte les configurations de gardiens à examiner, tandis que
la borne $O(|H|\,n\,3^n)$ tient également compte du coût de leur évaluation
sur la fenêtre.

Pour l'application numérique visée, $n=4$. Le nombre total de couples
candidats $(S,G)$ est donc $3^4-2^4=65$.
L'énumération exhaustive est ainsi négligeable et présente l'avantage
d'être exacte, transparente et indépendante des tolérances d'un solveur
MILP.
\end{remarque}

\subsection{Politique horairement reconfigurable}

Dans la politique reconfigurable définie par
\eqref{eq:cout_horaire_borne}, le choix des gardiens se décompose par heure.
La Proposition~\ref{pro:opt} et l'énumération précédente sont donc
appliquées indépendamment à chaque \(h\in H\). Cette politique peut changer
de gardiens au sein de \(H\), contrairement à la politique persistante qui
définit \(C_H^*\).

\subsection{Sous-additivité du coût opérationnel minimal}

La propriété structurelle centrale du coût optimal est maintenant établie.
Elle compare directement le fonctionnement d'une coalition réunie à celui de
deux coalitions séparées.

Les fonctions \(C_H^*\) et \(C_{H,\mathrm{hor}}^*\), définies sur
\(\mathcal S\), sont prolongées à \(2^N\) par la convention
\[
    C_H^*(\varnothing)
    =C_{H,\mathrm{hor}}^*(\varnothing)
    :=0.
\]

\begin{definition}[Sous-additivité]
Soit $f:2^N\to\mathbb R$ telle que $f(\varnothing)=0$. La fonction $f$ est dite
\textbf{sous-additive} si, pour tous \(A,B\subseteq N\) disjoints,
\[
    f(A\cup B)\leqslant f(A)+f(B).
\]
\end{definition}

\begin{theorem}[Sous-additivité des coûts sur la fenêtre]
\label{thm:sousadd}
Les fonctions \(C_H^*\) et \(C_{H,\mathrm{hor}}^*\) sont sous-additives.
\end{theorem}

\begin{proof}
Soient \(A,B\subseteq N\) disjoints. Si \(A=\varnothing\) ou
\(B=\varnothing\), l'inégalité résulte immédiatement de la convention
précédente. Supposons désormais \(A,B\in\mathcal S\).

Pour la politique persistante, posons
\(G:=G_H^*(A)\cup G_H^*(B)\). À chaque heure,
\(\mathbf t^h:=\mathbf t^{h,*}(A)+\mathbf t^{h,*}(B)\) sert
\(d^h(A\cup B)\). Les coalitions étant disjointes, les ensembles de gardiens
et les supports des répartitions le sont également. Ainsi,
\(G\in\mathcal G_H(A\cup B)\) et la solution construite a pour coût
\(C_H^*(A)+C_H^*(B)\). Par minimalité,
\[
    C_H^*(A\cup B)
    \leqslant C_H^*(A)+C_H^*(B).
\]

Pour la politique horairement reconfigurable, le même argument est appliqué
indépendamment à chaque heure, puis sommé sur \(H\). Il donne
\[
    C_{H,\mathrm{hor}}^*(A\cup B)
    \leqslant
    C_{H,\mathrm{hor}}^*(A)+C_{H,\mathrm{hor}}^*(B).
\]
\end{proof}

\begin{remarque}[La sous-additivité peut être non stricte]
Supposons que chaque opérateur sature sa capacité à chaque heure, c'est-à-dire
que \(d_i^h=q_i\) pour tous \(i\in N\) et \(h\in H\). Pour tout
\(S\in\mathcal S\), on a alors \(d^h(S)=q(S)\), de sorte qu'aucun
sous-ensemble strict de \(S\) ne peut servir sa demande. Par conséquent,
\[
    C_H^*(S)=C_{H,\mathrm{hor}}^*(S)
    =|H|\sum_{i\in S}\left(F_i+\gamma_iq_i\right).
\]
Les deux inégalités du Théorème~\ref{thm:sousadd} sont donc des égalités pour
toutes les coalitions disjointes. Leur réunion ne crée alors aucune économie
supplémentaire.
\end{remarque}

La sous-additivité stricte requiert ainsi une hypothèse supplémentaire. La
condition locale suivante isole un mécanisme suffisant : un seul des gardiens
retenus séparément peut acheminer toute la demande réunie sur la fenêtre.

\begin{proposition}[Condition locale de gain strict]
\label{prop:sousadd_stricte_locale}
Soient \(A,B\in\mathcal S\) deux coalitions disjointes. Supposons qu'il existe
\[
    \alpha\in
    \operatorname*{arg\,min}_{i\in G_H^*(A)\cup G_H^*(B)}\gamma_i
    \qquad\text{tel que}\qquad
    d^h(A\cup B)\leqslant q_\alpha,
    \quad \forall h\in H.
\]
Alors
\[
    C_H^*(A)+C_H^*(B)-C_H^*(A\cup B)
    \geqslant
    |H|\sum_{i\in\left(G_H^*(A)\cup G_H^*(B)\right)
    \setminus\{\alpha\}}F_i
    >0.
\]
En particulier, \(C_H^*(A\cup B)<C_H^*(A)+C_H^*(B)\).
\end{proposition}

\begin{proof}
Dans les solutions optimales séparées, le coût variable total est au moins
\(\gamma_\alpha D(A\cup B)\), puisque \(\gamma_\alpha\) est le plus faible
coût variable unitaire parmi les gardiens utilisés. Sous l'hypothèse de
capacité, \(\alpha\) peut servir seul \(A\cup B\) à chaque heure. Cette
solution admissible coûte
\[
    |H|F_\alpha+\gamma_\alpha D(A\cup B).
\]
La différence entre le coût des solutions séparées et celui de cette solution
donne la borne annoncée. Enfin, les deux coalitions étant disjointes, leurs
ensembles de gardiens le sont également et contiennent chacun au moins un
élément ; la stricte positivité découle donc de \(F_i>0\) pour tout \(i\in N\).
\end{proof}

\begin{corollaire}[Sous-additivité stricte en régime de faible trafic]
\label{cor:sousadd_stricte_faible_trafic}
Supposons que
\[
    d^h(N)\leqslant\min_{i\in N}q_i,
    \qquad \forall h\in H.
\]
Alors, pour toutes les coalitions disjointes \(A,B\in\mathcal S\),
\[
    C_H^*(A\cup B)<C_H^*(A)+C_H^*(B)
\]
et
\[
    C_{H,\mathrm{hor}}^*(A\cup B)
    <C_{H,\mathrm{hor}}^*(A)+C_{H,\mathrm{hor}}^*(B).
\]
\end{corollaire}

\begin{proof}
Pour la politique persistante, tout opérateur \(\alpha\in N\) vérifie
\(d^h(A\cup B)\leqslant d^h(N)\leqslant q_\alpha\) à chaque heure. La
Proposition~\ref{prop:sousadd_stricte_locale} s'applique donc au gardien de
plus faible \(\gamma_i\) parmi ceux des solutions optimales séparées. Pour la
politique horairement reconfigurable, le même raisonnement est appliqué à
chaque heure aux deux solutions horaires optimales. Il fournit une inégalité
stricte à chaque heure, dont la somme sur \(H\) reste stricte.
\end{proof}

\begin{remarque}[Portée et interprétation économique]
La sous-additivité garantit en général, sous chacune des deux politiques, que
la réunion de coalitions disjointes n'augmente pas leur coût opérationnel
total. L'exemple de saturation montre que ce gain peut être nul, tandis que le
Corollaire~\ref{cor:sousadd_stricte_faible_trafic} garantit un gain strict
lorsque chaque opérateur peut servir seul toute la demande. Ces propriétés
établissent l'intérêt collectif de la coopération, mais ne garantissent pas
que l'économie obtenue puisse être répartie de manière stable contre les
déviations de sous-coalitions. Comme le coût est une valorisation linéaire
positive de l'énergie consommée, la même conclusion vaut pour la consommation
énergétique totale.
\end{remarque}

\begin{corollaire}[Comparaison avec le fonctionnement autonome]
\label{cor:comparaison_autonome}
Pour tout $S\in\mathcal S$,
\[
C_H^*(S)
\leqslant
\sum_{i\in S}C_H^*(\{i\}),
\qquad
C_{H,\mathrm{hor}}^*(S)
\leqslant
\sum_{i\in S}C_H^*(\{i\}).
\]
\end{corollaire}

\begin{proof}
Il suffit d'appliquer successivement la sous-additivité de chaque fonction de
coût aux singletons disjoints qui composent $S$, puis d'utiliser
\(C_{H,\mathrm{hor}}^*(\{i\})=C_H^*(\{i\})\).
\end{proof}

\subsection{Superadditivité du jeu des économies}

La sous-additivité du coût doit être traduite dans le jeu défini à la
Sous-section~\ref{subsec:modele_jeu_economies}. La superadditivité obtenue
compare l'économie de la coalition réunie à la somme des économies que ses
membres pourraient réaliser en restant séparés.

\begin{proposition}[Propriétés des jeux sur la fenêtre]
\label{prop:proprietes_jeu_economies}
\label{prop:proprietes_jeu_fenetre}
Pour \(w\in\{v_H,v_{H,\mathrm{hor}}\}\),
\[
w(\varnothing)=0,
\qquad
w(\{i\})=0,
\qquad \forall i\in N,
\]
et
\[
w(S)\geqslant0,
\qquad \forall S\in\mathcal S.
\]
En outre, les deux jeux sont superadditifs : pour tous \(A,B\subseteq N\)
disjoints,
\[
w(A\cup B)\geqslant w(A)+w(B).
\]
Sous l'hypothèse du
Corollaire~\ref{cor:sousadd_stricte_faible_trafic}, cette inégalité est
stricte pour toutes les coalitions disjointes \(A,B\in\mathcal S\).
\end{proposition}

\begin{proof}
Les identités sur les singletons découlent de la référence autonome commune.
La non-négativité résulte du
Corollaire~\ref{cor:comparaison_autonome}. Pour chacun des deux jeux, la
superadditivité s'obtient en substituant la sous-additivité du coût
correspondant dans sa définition. La stricte sous-additivité établie au
Corollaire~\ref{cor:sousadd_stricte_faible_trafic} donne de même la stricte
superadditivité lorsque \(A,B\in\mathcal S\).
\end{proof}

\begin{corollaire}[Efficacité globale de la grande coalition]
\label{cor:efficacite_grande_coalition}
Pour toute partition \(\mathcal P\) de \(N\),
\[
    v_H(N)\geqslant\sum_{S\in\mathcal P}v_H(S),
    \qquad
    v_{H,\mathrm{hor}}(N)
    \geqslant
    \sum_{S\in\mathcal P}v_{H,\mathrm{hor}}(S).
\]
Sous l'hypothèse du
Corollaire~\ref{cor:sousadd_stricte_faible_trafic}, ces deux inégalités sont
strictes dès que \(\mathcal P\neq\{N\}\).
\end{corollaire}

\begin{proof}
Il suffit d'appliquer successivement la superadditivité de chaque jeu aux
coalitions disjointes de la partition \(\mathcal P\). Sous l'hypothèse de
faible trafic, la stricte superadditivité établie dans la
Proposition~\ref{prop:proprietes_jeu_economies} implique que toute fusion
successive des blocs d'une partition non triviale produit une inégalité
stricte.
\end{proof}

Ce corollaire donne le sens précis de l'efficacité collective : parmi toutes
les partitions des opérateurs, la grande coalition maximise l'économie totale.
L'inégalité est stricte sous la condition de faible trafic précédente, mais
l'efficacité, même stricte, ne garantit pas que l'économie puisse être répartie
de manière à satisfaire simultanément toutes les sous-coalitions. Cette
question constitue l'objet de la section suivante.

% \subsection{Cas particulier des capacités égales}

% Dans le cas particulier où les capacités sont égales, le problème de sélection optimale des gardiens se résout en temps polynomial, comme le montre la proposition suivante.

% \begin{proposition}[Cas de capacités égales]
% \label{prop:capacites_egales}
% Supposons que $\forall i\in S$, $q_i=q$.
% Alors le problème de sélection optimale des gardiens peut être résolu en temps $O(s^2\log s)$, où $s=|S|$.
% \end{proposition}

% \begin{proof}
% Posons
% \[
% k:=\left\lceil\frac{d(S)}{q}\right\rceil,
% \qquad
% \delta:=d(S)-(k-1)q\in(0,q].
% \]
% La faisabilité impose au moins $k$ gardiens. Réciproquement, la
% répartition gloutonne n'achemine du trafic que sur $k$ équipements lorsque
% les capacités valent toutes $q$. Si une solution activait davantage de
% gardiens, on pourrait choisir une répartition gloutonne de même coût
% variable, puis placer en veille tout équipement n'acheminant aucun trafic,
% ce qui réduirait strictement le coût fixe. Toute solution optimale active
% donc exactement $\lceil d(S)/q\rceil$ gardiens.

% Dans un ordre glouton, les $k-1$ premiers gardiens sont saturés et le
% pivot achemine $\delta$. Pour chaque opérateur $j\in S$ candidat au rôle
% de pivot, posons
% \[
% P_j:=\{i\in S\setminus\{j\}\mid\gamma_i\leqslant\gamma_j\}.
% \]
% Le candidat est admissible si $|P_j|\geqslant k-1$. La meilleure solution
% ayant $j$ pour pivot choisit alors dans $P_j$ les $k-1$ opérateurs de plus
% faible valeur $F_i+\gamma_iq$. Si $A_j$ désigne cet ensemble, son coût est
% \[
% F_j+\gamma_j\delta
% +\sum_{i\in A_j}\left(F_i+\gamma_iq\right).
% \]
% Pour chacun des $s$ pivots candidats, la sélection de $A_j$ peut être
% effectuée par un tri en $O(s\log s)$. Tous les candidats sont donc examinés
% en $O(s^2\log s)$.
% \end{proof}

% En répétant cette optimisation pour chaque coalition, nous obtenons le coût de
% la coopération globale ainsi que ceux des accords que pourraient former des
% groupes plus petits. La section suivante transforme ces coûts en économies de
% coopération et les utilise pour étudier l'intérêt collectif et la stabilité de
% l'accord global.
